% !TeX root = proyecto.tex

%=========================================================
\chapter{Introducción}

En las unidades del IPN, la alta demanda de alimentos durante los breves periodos de receso genera ineficiencias, largas filas y desanima a la comunidad politécnica a consumir sus alimentos en las limitadas cafeterías de sus unidades.

El ejemplo más cercano es la Escuela Superior de Cómputo (ESCOM) \cite{ipn2025estadistica} donde aproximadamente 2,842 alumnos del turno matutino deben ser atendidos en recesos de apenas 30 minutos por solo una cafetería y una barra de café que manejan más de 100 pedidos en el turno.

Un fenómeno similar se abarca en el estudio de León Martínez en \cite{leon2023planificacion}, el cual muestra que la dosificación de horarios de receso y consumo de alimentos tiene complicaciones propias agravadas por una capacidad limitada de los centros de alimentos incluidos en los planteles ante una comunidad que incluye a estudiantes, docentes y personal administrativo.

Los sistemas tradicionales de gestión de pedidos basados en papel o pizarrón y las soluciones digitales comerciales existentes no se adaptan a las necesidades específicas de la logística académica de alta rotación \cite{lavu2025kds}, por ejemplo, los casos estudiados por García Martínez, R. et al. \cite{garcia2024estrategia} y por Gelvez Jurado, A. H., et al. \cite{gelvez2024estudio} demuestran cómo las herramientas analógicas pueden causar demoras y conflictos en la línea de producción de alimentos preparados en cafeterías escolares.

Ante esto, el presente proyecto propone el desarrollo de un Sistema de Gestión de Cocina (KMS, por sus siglas en inglés) como aplicación web, diseñada específicamente para las condiciones de las 3 cafeterías que prestan sus servicios en ESIME Zacatenco y ESFM y gestione correctamente el flujo de trabajo, la disponibilidad de insumos y la coordinación entre las estaciones de trabajo en la cocina y con las demás estaciones de carácter operativo \cite{zaenal2024commercial}. Para comprender mejor la necesidad que se trata de cubrir, a continuación, se abordan los antecedentes.

\section{Presentación}


\cdtInstrucciones{
	Presente en un par de párrafos el contexto y el problema en que se define el proytecto.
}

\cdtInstrucciones{
	Indique el propósito del documento, a quien va dirigido y como debe ser utilizado el documento.
}
	
%---------------------------------------------------------
\section{Organización del contenido}

\cdtInstrucciones{
	Indique el contenido y organización del documento.
}
	En el capítulo \ref{cap:reqUsr} ...
	
	En el capítulo \ref{cap:reqSist} ...

%---------------------------------------------------------
\section{Notación, símbolos y convenciones utilizadas}

\cdtInstrucciones{
	Indique la notacion utilizada así como nórmas o estándares de documentación utilizados en el documento.
}	

